\documentclass[a4paper,10pt]{report}\input{../0.Préambule/Préambule_cours_prof}\begin{document}

\newpage
\section*{Théorème de Pythagore}
\addcontentsline{toc}{section}{Théorème de Pythagore}

\begin{exer1}
Pour chaque triangle, justifier pourquoi on a l'égalité de Pythagore et énoncer cette propriété.
\begin{enumerate}[font=\color{exer1}]
\begin{tabularx}{\linewidth}{*{2}{>{\centering \arraybackslash}X}}
\item \definecolor{qqwuqq}{rgb}{0.,0.39215686274509803,0.}
\begin{tikzpicture}[line cap=round,line join=round,>=triangle 45,x=1.0cm,y=1.0cm,scale=1]
\clip(0.5,2.5) rectangle (5.5,6.5);
\draw[line width=1.2pt,color=qqwuqq,fill=qqwuqq,fill opacity=0.1] (1.,5.575735931288071) -- (1.4242640687119288,5.575735931288071) -- (1.4242640687119288,6.) -- (1.,6.) -- cycle; 
\draw [line width=1.2pt] (1.,6.)-- (1.,3.);
\draw [line width=1.2pt] (1.,3.)-- (5.,6.);
\draw [line width=1.2pt] (5.,6.)-- (1.,6.);
\begin{scriptsize}
\draw [color=blue] (1.,6.)-- ++(-3.5pt,0 pt) -- ++(7.0pt,0 pt) ++(-3.5pt,-3.5pt) -- ++(0 pt,7.0pt);
\draw[color=blue] (0.75,6.145) node {\large{$A$}};
\draw [color=blue] (1.,3.)-- ++(-3.5pt,0 pt) -- ++(7.0pt,0 pt) ++(-3.5pt,-3.5pt) -- ++(0 pt,7.0pt);
\draw[color=blue] (0.75,2.7) node {\large{$B$}};
\draw [color=blue] (5.,6.)-- ++(-3.5pt,0 pt) -- ++(7.0pt,0 pt) ++(-3.5pt,-3.5pt) -- ++(0 pt,7.0pt);
\draw[color=blue] (5.138,6.2) node {\large{$C$}};
\draw [color=blue] (6.,3.)-- ++(-3.5pt,0 pt) -- ++(7.0pt,0 pt) ++(-3.5pt,-3.5pt) -- ++(0 pt,7.0pt);
\end{scriptsize}
\end{tikzpicture}&
\item \definecolor{qqwuqq}{rgb}{0.,0.39215686274509803,0.}
\begin{tikzpicture}[line cap=round,line join=round,>=triangle 45,x=1.0cm,y=1.0cm,scale=1]
\clip(5.5,2.5) rectangle (12.5,6.5);
\draw [shift={(6.,3.)},line width=1.2pt,color=qqwuqq,fill=qqwuqq,fill opacity=0.1] (0,0) -- (0.:0.6) arc (0.:36.982008980943:0.6) -- cycle;
\draw [shift={(12.,3.)},line width=1.2pt,color=qqwuqq,fill=qqwuqq,fill opacity=0.1] (0,0) -- (126.982008980943:0.6) arc (126.982008980943:180.:0.6) -- cycle;

\draw [line width=1.2pt] (6.,3.)-- (9.828722935828715,5.883265214237399);
\draw [line width=1.2pt] (9.828722935828715,5.883265214237399)-- (12.,3.);
\draw [line width=1.2pt] (6.,3.)-- (12.,3.);

\begin{scriptsize}
\draw [color=blue] (6.,3.)-- ++(-3.5pt,0 pt) -- ++(7.0pt,0 pt) ++(-3.5pt,-3.5pt) -- ++(0 pt,7.0pt);
\draw[color=blue] (6.138,3.45) node {\large{$E$}};
\draw [color=blue] (12.,3.)-- ++(-3.5pt,0 pt) -- ++(7.0pt,0 pt) ++(-3.5pt,-3.5pt) -- ++(0 pt,7.0pt);
\draw[color=blue] (12.138,3.45) node {\large{$F$}};
\draw [color=blue] (9.8,5.8)-- ++(-3.5pt,0 pt) -- ++(7.0pt,0 pt) ++(-3.5pt,-3.5pt) -- ++(0 pt,7.0pt);
\draw[color=blue] (9.978,6.25) node {\large{$D$}};
\draw [color=blue] (1.,-5.)-- ++(-3.5pt,0 pt) -- ++(7.0pt,0 pt) ++(-3.5pt,-3.5pt) -- ++(0 pt,7.0pt);
\draw[color=qqwuqq] (7.778,3.45) node {\large{$37\textrm{\degre}$}};
\draw[color=qqwuqq] (10.958,3.53) node {\large{$53\textrm{\degre}$}};
\end{scriptsize}
\end{tikzpicture}\\
\item \definecolor{qqwuqq}{rgb}{0.,0.39215686274509803,0.}
\begin{tikzpicture}[line cap=round,line join=round,>=triangle 45,x=1.0cm,y=1.0cm,scale=1]
\clip(0.5,-5.5) rectangle (7.5,-1.5);
\draw [shift={(1.,-5.)},line width=1.2pt,color=qqwuqq,fill=qqwuqq,fill opacity=0.1] (0,0) -- (0.:0.6) arc (0.:45.:0.6) -- cycle;
\draw [shift={(7.,-5.)},line width=1.2pt,color=qqwuqq,fill=qqwuqq,fill opacity=0.1] (0,0) -- (135.:0.6) arc (135.:180.:0.6) -- cycle;
\draw [line width=1.2pt] (1.,-5.)-- (4.,-2.);
\draw [line width=1.2pt] (2.401005050633883,-3.457573593128807) -- (2.5424264068711926,-3.5989949493661175);
\draw [line width=1.2pt] (2.457573593128807,-3.4010050506338825) -- (2.5989949493661166,-3.542426406871193);
\draw [line width=1.2pt] (4.,-2.)-- (7.,-5.);
\draw [line width=1.2pt] (5.542426406871193,-3.4010050506338825) -- (5.4010050506338825,-3.542426406871193);
\draw [line width=1.2pt] (5.5989949493661175,-3.457573593128807) -- (5.4575735931288065,-3.5989949493661175);
\draw [line width=1.2pt] (7.,-5.)-- (1.,-5.);
\begin{scriptsize}
\draw [color=blue] (1.,-5.)-- ++(-3.5pt,0 pt) -- ++(7.0pt,0 pt) ++(-3.5pt,-3.5pt) -- ++(0 pt,7.0pt);
\draw[color=blue] (1,-4.55) node {\large{$G$}};
\draw [color=blue] (7.,-5.)-- ++(-3.5pt,0 pt) -- ++(7.0pt,0 pt) ++(-3.5pt,-3.5pt) -- ++(0 pt,7.0pt);
\draw[color=blue] (7.138,-4.55) node {\large{$H$}};
\draw [color=blue] (4.,-2.)-- ++(-3.5pt,0 pt) -- ++(7.0pt,0 pt) ++(-3.5pt,-3.5pt) -- ++(0 pt,7.0pt);
\draw[color=blue] (4.138,-1.75) node {\large{$U$}};
\draw [color=blue] (8.,-5.)-- ++(-3.5pt,0 pt) -- ++(7.0pt,0 pt) ++(-3.5pt,-3.5pt) -- ++(0 pt,7.0pt);
\draw[color=qqwuqq] (2.538,-4.55) node {\large{$45\textrm{\degre}$}};
\draw[color=qqwuqq] (5.998,-4.61) node {\large{$45\textrm{\degre}$}};
\end{scriptsize}
\end{tikzpicture}&
\item \definecolor{qqwuqq}{rgb}{0.,0.39215686274509803,0.}
\begin{tikzpicture}[line cap=round,line join=round,>=triangle 45,x=1.0cm,y=1.0cm,scale=1]
\clip(7.5,-5.5) rectangle (13.5,-1.5);
\draw[line width=1.2pt,color=qqwuqq,fill=qqwuqq,fill opacity=0.1] (13.,-4.575735931288071) -- (12.57573593128807,-4.575735931288071) -- (12.57573593128807,-5.) -- (13.,-5.) -- cycle; 
\draw[line width=1.2pt,color=qqwuqq,fill=qqwuqq,fill opacity=0.1] (8.,-2.4242640687119286) -- (8.42426406871193,-2.4242640687119286) -- (8.42426406871193,-2.) -- (8.,-2.) -- cycle; 
\draw [line width=1.2pt] (8.,-5.)-- (8.,-2.);
\draw [line width=1.2pt] (8.,-2.)-- (13.,-2.);
\draw [line width=1.2pt] (13.,-2.)-- (13.,-5.);
\draw [line width=1.2pt] (13.,-5.)-- (8.,-5.);
\draw [line width=1.2pt] (8.,-5.)-- (13.,-2.);
\begin{scriptsize}
\draw [color=blue] (8.,-5.)-- ++(-3.5pt,0 pt) -- ++(7.0pt,0 pt) ++(-3.5pt,-3.5pt) -- ++(0 pt,7.0pt);
\draw[color=blue] (7.75,-4.55) node {\large{$J$}};
\draw [color=blue] (8.,-2.)-- ++(-3.5pt,0 pt) -- ++(7.0pt,0 pt) ++(-3.5pt,-3.5pt) -- ++(0 pt,7.0pt);
\draw[color=blue] (8.138,-1.75) node {\large{$K$}};
\draw [color=blue] (13.,-2.)-- ++(-3.5pt,0 pt) -- ++(7.0pt,0 pt) ++(-3.5pt,-3.5pt) -- ++(0 pt,7.0pt);
\draw[color=blue] (13.138,-1.75) node {\large{$L$}};
\draw [color=blue] (13.,-5.)-- ++(-3.5pt,0 pt) -- ++(7.0pt,0 pt) ++(-3.5pt,-3.5pt) -- ++(0 pt,7.0pt);
\draw[color=blue] (13.25,-4.55) node {\large{$M$}};
\end{scriptsize}
\end{tikzpicture}\\
\end{tabularx}
\end{enumerate}
\end{exer1}

\begin{exer1}
Dans chacun des cas, nommer l'hypoténuse des triangles.
\begin{enumerate}[font=\color{exer1}]
\item $ART$ est rectangle en $A$.
\item $[BU]$ et $[SB]$ sont les côtés de l'angle droit.
\end{enumerate}
\end{exer1}

\begin{exer1}

\vspace{0.5cm}
\noindent Nommer tous les triangles rectangles de la \\figure et, pour chacun, nommer son hypoténuse.

\vspace{-3cm}
\begin{flushright}
\definecolor{qqwuqq}{rgb}{0.,0.39215686274509803,0.}
\begin{tikzpicture}[line cap=round,line join=round,>=triangle 45,x=1.0cm,y=1.0cm,scale=0.9]
\clip(0.4,0.4) rectangle (7.7,5.7);
\draw[line width=1.2pt,color=qqwuqq,fill=qqwuqq,fill opacity=0.1] (1.2932875073293875,3.943027044276394) -- (1.7154604282114296,3.900955380725676) -- (1.7575320917621478,4.323128301607718) -- (1.3353591708801056,4.365199965158436) -- cycle; 
\draw[line width=1.2pt,color=qqwuqq,fill=qqwuqq,fill opacity=0.1] (2.810263340389897,4.620526680779794) -- (3.189736659610103,4.430790021169692) -- (3.3794733192202053,4.810263340389898) -- (3.,5.) -- cycle; 
\draw[line width=1.2pt,color=qqwuqq,fill=qqwuqq,fill opacity=0.1] (4.6605887450304575,3.745441558772843) -- (4.915147186257614,3.4060303038033) -- (5.254558441227157,3.6605887450304575) -- (5.,4.) -- cycle; 
\draw[line width=1.2pt,color=qqwuqq,fill=qqwuqq,fill opacity=0.1] (5.424264068711929,1.) -- (5.424264068711929,1.4242640687119283) -- (5.,1.4242640687119283) -- (5.,1.) -- cycle; 
\draw [line width=1.2pt] (1.,1.)-- (5.,4.);
\draw [line width=1.2pt] (1.,1.)-- (5.,1.);
\draw [line width=1.2pt] (5.,1.)-- (5.,4.);
\draw [line width=1.2pt] (5.,4.)-- (3.,5.);
\draw [line width=1.2pt] (3.,5.)-- (1.,1.);
\draw [line width=1.2pt] (1.,1.)-- (1.3353591708801056,4.365199965158436);
\draw [line width=1.2pt] (1.3353591708801056,4.365199965158436)-- (5.,4.);
\draw [line width=1.2pt] (1.,1.)-- (7.25,1.);
\draw [line width=1.2pt] (7.25,1.)-- (5.,4.);
\begin{scriptsize}
\draw [color=blue] (1.,1.)-- ++(-3.5pt,0 pt) -- ++(7.0pt,0 pt) ++(-3.5pt,-3.5pt) -- ++(0 pt,7.0pt);
\draw[color=blue] (0.7,0.63) node {\large{$A$}};
\draw [color=blue] (5.,4.)-- ++(-3.5pt,0 pt) -- ++(7.0pt,0 pt) ++(-3.5pt,-3.5pt) -- ++(0 pt,7.0pt);
\draw[color=blue] (5.44,4.55) node {\large{$B$}};
\draw [color=blue] (3.,5.)-- ++(-3.5pt,0 pt) -- ++(7.0pt,0 pt) ++(-3.5pt,-3.5pt) -- ++(0 pt,7.0pt);
\draw[color=blue] (3.36,5.4) node {\large{$C$}};
\draw [color=blue] (1.3353591708801056,4.365199965158436)-- ++(-3.5pt,0 pt) -- ++(7.0pt,0 pt) ++(-3.5pt,-3.5pt) -- ++(0 pt,7.0pt);
\draw[color=blue] (1.48,4.95) node {\large{$D$}};
\draw [color=blue] (5.,1.)-- ++(-3.5pt,0 pt) -- ++(7.0pt,0 pt) ++(-3.5pt,-3.5pt) -- ++(0 pt,7.0pt);
\draw[color=blue] (4.98,0.65) node {\large{$E$}};
\draw [color=blue] (7.25,1.)-- ++(-3.5pt,0 pt) -- ++(7.0pt,0 pt) ++(-3.5pt,-3.5pt) -- ++(0 pt,7.0pt);
\draw[color=blue] (7.48,0.63) node {\large{$F$}};
\end{scriptsize}
\end{tikzpicture}
\end{flushright}
\end{exer1}

\begin{exer1}
Le triangle $GAL$, rectangle en $A$, est tel que $GA=84$m et $AL= 35$m.

\medskip
\noindent Calculer la longueur $GL$ de son hypoténuse.
\end{exer1}

\begin{exer1}
A quel triangle s'applique chacune des égalités suivantes ?
\begin{enumerate}[font=\color{exer1}]
\item $AB^2=AC^2+BC^2$              
\item $BC^2=AB^2+AC^2$                
\item $AC^2=AB^2+BC^2$
\end{enumerate}
\end{exer1}

\begin{exer1}
Le triangle $PIM$, rectangle en $P$, est tel que $PI=68$mm et $MI=68,9$mm.

\medskip
\noindent Calculer la longueur du côté $[PM]$
\end{exer1}

\begin{exer1}
$BUS$ est un triangle rectangle en $B$ tel que $BS= 6$cm et $SU=9$cm

\medskip
\noindent Calculer la longueur $UB$, arrondie au millimètre près.
\end{exer1}

\begin{exer1}
$ABC$ est un triangle rectangle en $A$ tel que $AB= 3$cm et $AC= 1$cm.
\begin{enumerate}[font=\color{exer1}]
\item Joseph a écrit \og $BC^2=6+2$ ;$ BC^2=8$ donc $BC=4$cm \fg{}
Indique et analyse ses erreurs.
\item Calcule $BC^2$ puis en utilisant la touche racine carrée de la calculatrice,\\ donne la valeur de $BC$ approchée par défaut au millimètre près.
\end{enumerate}
\end{exer1}

\begin{exer1}
$MER$ est un triangle rectangle en $E$.

\vspace{-3.5cm}
\begin{flushright}
\definecolor{qqwuqq}{rgb}{0.,0.39215686274509803,0.}
\begin{tikzpicture}[line cap=round,line join=round,>=triangle 45,x=1.0cm,y=1.0cm]
\clip(0.4,0.4) rectangle (6.4,4.2);
\draw[line width=1.2pt,color=qqwuqq,fill=qqwuqq,fill opacity=0.1] (4.137733595215663,3.144597732509787) -- (4.377135862705876,2.79433132772545) -- (4.727402267490212,3.033733595215662) -- (4.488,3.384) -- cycle; 
\draw [line width=1.2pt] (1.,1.)-- (5.68,1.64);
\draw [line width=1.2pt] (5.68,1.64)-- (4.488,3.384);
\draw [line width=1.2pt] (4.488,3.384)-- (1.,1.);
\begin{scriptsize}
\draw [color=blue] (1.,1.)-- ++(-3.5pt,0 pt) -- ++(7.0pt,0 pt) ++(-3.5pt,-3.5pt) -- ++(0 pt,7.0pt);
\draw[color=blue] (0.7,0.63) node {\large{$M$}};
\draw [color=blue] (5.68,1.64)-- ++(-3.5pt,0 pt) -- ++(7.0pt,0 pt) ++(-3.5pt,-3.5pt) -- ++(0 pt,7.0pt);
\draw[color=blue] (6.12,2.19) node {\large{$R$}};
\draw [color=blue] (4.488,3.384)-- ++(-3.5pt,0 pt) -- ++(7.0pt,0 pt) ++(-3.5pt,-3.5pt) -- ++(0 pt,7.0pt);
\draw[color=blue] (4.84,3.91) node {\large{$E$}};
\end{scriptsize}
\end{tikzpicture}
\end{flushright}

\vspace{-10mm}
\begin{enumerate}[font=\color{exer1}]
\item Écris l'égalité de Pythagore pour ce triangle.
\item Le tableau suivant présente plusieurs cas de dimensions du triangle $MER$. Complète-le en écrivant le détail de tes calculs (tu arrondiras au dixième si nécessaire) :

\renewcommand{\arraystretch}{1.9}
\begin{tabularx}{\linewidth}{|>{\columncolor{exer1!20}}p{1cm}|*{5}{>{\centering \arraybackslash}X|}}
\hline
\rowcolor{exer1!20} & Cas $1$ & Cas $2$ & Cas $3$ &Cas $4$ & Cas $5$\\\hline
$MR$ & \dotfill & \dotfill & $5,3$cm & $9,1$cm & $7$m\\\hline
$RE$ & $15$cm & $36$cm & \dotfill & $9$cm & \dotfill m\\\hline
$ME$ & $8$cm & $7,7$dm & $2,8$cm & \dotfill & $53$cm\\\hline
\end{tabularx}

\end{enumerate}
\end{exer1}

\begin{exer2}
Théo veut franchir, avec une échelle, un mur de $3,50$m de haut devant lequel se trouve un fossé rempli d'eau, d'une largeur de $1,15$m.
\begin{enumerate}[font=\color{exer2}]
\item Fais un schéma de la situation.
\item Il doit poser l'échelle sur le sommet du mur. Quelle doit être la longueur minimum de l'échelle ?
\end{enumerate}
\end{exer2}

\begin{exer2}
$TSF$ est un triangle isocèle en $S$ tel que $ST= 4,5$cm et $TF=5,4$cm.
\begin{enumerate}[font=\color{exer2}]
\item Calcule la longueur de la hauteur relative à la base $[TF]$.
\item Déduis-en l'aire de ce triangle.
\end{enumerate}
\end{exer2}

\begin{exer2}
Dans chacun des cas ci-dessous :
\begin{enumerate}[font=\color{exer2}]
\item[•]identifie le plus long côté du triangle $EFG$ ;
\item[•]calcule, d'une part, le carré de la longueur de ce côté ;
\item[•]calcule, d'autre part, la somme des carrés des longueurs des deux autres côtés ;
\item[•]compare les résultats obtenus et conclus.
\end{enumerate}
\begin{enumerate}[font=\color{exer2}]
\item $EF=4,5$cm ; $FG= 6$cm ;$ EG= 7,5$cm.
\item $EF= 3,6$cm ; $FG= 6$cm ; $EG= 7$cm.
\item $FG= 64$mm ; $EF= 72$mm ; $EG= 65$mm.
\item $EF= 3,2$dam ; $FG= 25,6$m ; $EG= 19,2$m.
\end{enumerate}
\end{exer2}

\begin{exer2}
Le triangle $OUI$ est tel que : $UI= 5$cm ; $UO=1,4$cm et $OI= 4,8$cm.
\begin{enumerate}[font=\color{exer2}]
\item Construis ce triangle en vraie grandeur.
\item Par la symétrie de centre $O$, construis les points $T$ et $N$ symétriques respectifs des points $U$ et $I$.
\item Quelle semble être la nature de $NUIT$ ? Démontre ta conjecture.
\end{enumerate}
\end{exer2}

\begin{exer2}
$ABC$ est un triangle rectangle en $B$ tel que : $AB=5$cm et $AC= 8$cm.
\begin{enumerate}[font=\color{exer2}]
\item Calcule $BC$ (arrondis au mm).
\item $D$ est un point tel que : $CD=20$cm et $BD= 19$cm. $D$ est-il unique ?
\item Montre que le triangle $BCD$ est rectangle. Précise en quel point.
\item Déduis-en que les points $A$, $B$ et $D$ sont alignés.
\end{enumerate}
\end{exer2}

\begin{exer3}
Deux droites $(d_1)$ et $(d_2)$ sont sécantes en $O$ ;
$M$ est un point de $(d_1)$ tel que : $OM=11,9$cm et $N$ est un point de $(d_2)$ tel que : $ON= 12$cm. On sait d'autre part que : $MN= 16,9$cm.


\medskip
\noindent Démontre que les droites $(d_1)$ et $(d_2)$ sont perpendiculaires.
\end{exer3}

\end{document}